\lecture{16}
\title{Sparse Representations: NNs vs. Kernels}

\begin{document}

\begin{remark}
    These notes are not particularly well-written,
    as they gloss over most of the computational details
    and only summarize the main conceptual points
    of this lecture, occasionally using mathematical informalities.
\end{remark}

\subsection{NNs vs. Kernels}

Recall from \href{/6.790/lectures/6/#kernels-unlimited-features}{Lecture 6}
and \href{/6.790/lectures/12}{Lecture 12}
that neural networks and kernel methods
are both linear models over 
the hidden layers of a neural network
encode the features that define the kernel.

The purpose of this lecture
is to roughly explain how neural networks
may be advantageous relative to kernel methods
when the model being learned
has a \emph{sparse representation}.

A model is said to have a \textbf{sparse representation}
if it is a map $\mathbb{R}^d \to \mathbb{R}$
for which only $r \ll d$ of the $d$ input dimensions
are relevant.
For example, $f \colon \mathbb{R}^{1000} \to \mathbb{R}$
via $f(x) := x_1 + x_1x_2 + x_1x_2x_3$ has a sparse representation,
with $r = 3$ and $d = 1000$.

\subsection{Kernel Performance}

In the $f \colon \mathbb{R}^{1000} \to \mathbb{R}$ toy example,
kernel methods will produce a hypothesis
$h(x) := w^{\top} \phi(x)$, where
$\phi(x)$ is a vector of dimension at least
$\binom{1000}{3} = O(1000^3)$.
Therefore, learning this hypothesis
requires $O(1000^3)$ samples.
(Justification ommitted.)

That's a lot of samples!  Can NNs do better?

\subsection{NN Performance}

The advantage of NNs is the following:
\begin{itemize}
    \item Kernel methods use SGD to directly attack the question,
    ``Which of the $O(1000^3)$ features are relevant?''
    \item A neural network ultimately uses SGD to answer the same question,
    but it does so incrementally.
    \begin{itemize}
        \item Its first layer uses SGD to answer the question,
        ``Which $r = 3$ of the $d = 1000$ input dimensions are relevant?''
        \item Then the deeper hidden layers need only
        use SGD to answer the question,
        ``Which of the $O(r^3)$ features are relevant?''
        \item The advantage for NNs is especially evident
        given \emph{sparse representations} (i.e., $r \ll d$).
    \end{itemize}
\end{itemize}

\textbf{Reminder.} The first hidden layer does \emph{not} encode
the same set of $O(1000^3)$ features that the kernel encodes!
A single neuron in the first hidden layer
can only represent $\mathrm{ReLU}(w^{\top}x)$, after all---not
even close to anything like $x_1x_2x_3$.
It is only the final output of \emph{all} of the hidden layers together
that is capable of learning all $O(1000^3)$ features.

\begin{remark}
    Technically, as
    \href{/6.790/lectures/12/#neural-network-universal-approximation}{Lecture 12}
    describes, a linear combination of
    the neurons in just a single hidden layer of an NN
    is capable of approximating any continuous function arbitrarily well,
    so such a combination could represent all $O(1000^3)$ features,
    so long as enough neurons are used.

    But in practice, using just a single layer is infeasible;
    recalling \href{/6.790/lectures/13/#neural-network-depth}{Lecture 13},
    the number of parameters in an NN can be reduced
    by increasing the number of hidden layers.
\end{remark}

Let's explain, now, how SGD allows NNs
to answer the first question---``Which $r = 3$
of the $d = 1000$ input dimensions are relevant?''---using
the weights $W \in \mathbb{R}^{m \times 1000}$
between its input layer and its width-$m$ first hidden layer.

The point is that each iteration of SGD
might shift any individual weight of $W$ by $O((1000m)^{-1})$,
rather than $O(1000^{-3})$,
because there are only $1000m$ parameters in $W$.
So here's what happens:
\begin{itemize}
    \item The term $x_1$ inside of $f(x)$
    incentivizes increasing the weights of $w_{j, 1}$
    within $O(1000m)$ SGD iterations.
    \item Now that $w_{j, 1}$ is large,
    the term $x_1x_2$ incentivizes
    increasing the weights of $w_{j, 2}$
    within another $O(1000m)$ SGD iterations.
    \item With $w_{j, 1}$ and $w_{j, 2}$ large,
    the term $x_1x_2x_3$ incentivizes increasing
    the weights of $w_{j, 3}$ after $O(1000m)$ more SGD iterations.
\end{itemize}
So it only takes $O(dmr)$ SGD iterations
to identify which $r$ input dimensions are relevant,
after which only $O(r^3)$ SGD steps are needed
to determine which of the $O(r^3)$ features over those
dimensions are relevant.

\begin{remark}
    In the first $O(1000m)$ SGD iterations,
    there is little incentive to increase the weights of
    $w_{j, 2}$ or $w_{j, 3}$.
    This is because $w_{j, 1}$ is not large enough---more precisely,
    the $\frac{\partial \mathcal{R}}{\partial w_{j, 2}}$ update
    is proportional to $w_{j, 1}$.

    Thus, NNs perform best for \textbf{staircase functions},
    whose terms each extend the previous by one coordinate;
    it is easier to learn $f(x) = x_1 + x_1x_2$ than $f(x) = x_1 + x_2x_3$.
\end{remark}

\end{document}